\section{Certificate formats and API contracts}
\label{app:formats}

\subsection{Canonical source objects}

A polynomial is a sorted array of triples
\code{[rate, degree, coefficient]}. Rates and coefficients are strings in the
canonical syntax of Python \code{Fraction}: an integer or a reduced fraction
with a positive denominator. Examples are \code{"-2"}, \code{"1/2"}, and
\code{"0"}. Inputs such as \code{"2/2"}, \code{"01"}, floats, and Boolean
values are rejected by the decoder; a Boolean must not pass merely because
Python treats it as an integer subtype. Zero coefficients, repeated keys, and
unsorted rows are rejected at the encoded boundary.

For example, $e^x-1-x$ has the source polynomial
\par\noindent\begin{minipage}{\linewidth}
\begin{lstlisting}
[["0", 0, "-1"], ["0", 1, "-1"], ["1", 0, "1"]]
\end{lstlisting}
\end{minipage}\par
and a source request can be
\par\noindent\begin{minipage}{\linewidth}
\begin{lstlisting}
{
  "polynomial": [["0",0,"-1"], ["0",1,"-1"], ["1",0,"1"]],
  "goal": {"kind": "positive_open_ray", "anchor": "0"}
}
\end{lstlisting}
\end{minipage}\par

These JSON objects are the prototype's independently supplied source problems.
They are not a serialization of arbitrary Lean expressions. A Lean implementation
must retain the separate source-evaluation equality discussed in Section
\ref{sec:lean}.

\subsection{Ladder receipts}

\par\noindent\begin{minipage}{\linewidth}
\begin{lstlisting}
{
  "kind": "ladder-v1",
  "anchor": "0",
  "bits": 96,
  "order": 32,
  "steps": [
    {
      "rho": "0",
      "polynomial": [["0",0,"-1"], ["0",1,"-1"], ["1",0,"1"]],
      "lower": "0", "upper": "0"
    },
    {
      "rho": "0",
      "polynomial": [["0",0,"-1"], ["1",0,"1"]],
      "lower": "0", "upper": "0"
    },
    {
      "rho": "1", "polynomial": [["1",0,"1"]],
      "lower": "1", "upper": "1"
    }
  ]
}
\end{lstlisting}
\end{minipage}\par

The terminal zero is implicit and must be computed after the final factor.
The precision parameters select a reproducible enclosure calculation, not a
numerical trust tolerance. At a nonzero anchor the stored endpoints concern
the positively normalised value of equation~\eqref{eq:normalised-evaluation}.

\subsection{Tail receipts}

\par\noindent\begin{minipage}{\linewidth}
\begin{lstlisting}
{
  "kind": "tail-v1",
  "sign": 1,
  "cutoff": "2880000",
  "orders": [["0", 6]]
}
\end{lstlisting}
\end{minipage}\par

For the separately supplied input $e^x-1000x^5$, this is the certificate from
Theorem~\ref{thm:tail}. The checker extracts the leading rate, degree, coefficient,
and all block norms from the source; it does not accept those facts merely
because the receipt states them. The zero source admits a separate sign-zero
receipt with no lower-rate orders.

\subsection{The accepted source request types}

\begin{center}
\begin{tabularx}{\linewidth}{@{}p{0.34\linewidth}X@{}}
\toprule
Request & Required receipt conclusion\\
\midrule
\code{nonnegative_ray} & A ladder with the identical anchor\\
\code{positive_open_ray} & Identical anchor and at least one strictly positive initial lower bound\\
\code{positive_closed_ray} & Identical anchor and strictly positive first initial lower bound\\
\code{eventual_sign} & A tail receipt with the requested sign; cutoff is an explicit witness\\
\code{nonnegative_interval} & A complete cover of exactly the requested interval\\
\code{counterexample_interval} & A certified negative point in exactly that interval\\
\bottomrule
\end{tabularx}
\end{center}

The public functions \code{verify(sourcePolynomial, receipt)} and
\code{verify_goal(sourceProblem, receipt)} answer different questions. The
former establishes the receipt's own conclusion against an independent
polynomial; the latter also checks that this conclusion answers the actual
request. Code invoking only the first function has not implemented the complete
source-goal contract.
